% SPDX-FileCopyrightText: 2025 IObundle
%
% SPDX-License-Identifier: MIT

Before anything else, software must call \texttt{versat\_init}, passing the
accelerator's base address, to initialize the generated driver. Once
initialized, a run is started and stopped through the generated API
described in Section~\ref{sec:if_config_state}.

This API is declared in a generated \texttt{versat\_accel.h}, produced by
instantiating the template \texttt{software/\allowbreak templates/\allowbreak versat\_header.tpl}
(via \texttt{software/\allowbreak common/\allowbreak templateEngine.cpp}): every \texttt{@\{...\}}
placeholder below is substituted at code-generation time, for example
\texttt{@\{typeName\}} with the top-level module's name.

\lstinputlisting[language=C,firstline=49,lastline=70,caption={\texttt{software/templates/versat\_header.tpl} (lines 49--70): run-control API declared in every generated \texttt{versat\_accel.h}.},label={lst:versat_header}]{../software/templates/versat_header.tpl}

Each generated header also declares two pointers, substituting \texttt{@\{typeName\}} for the top-level module's name:

\lstinputlisting[language=C,firstline=147,lastline=148,caption={\texttt{software/templates/versat\_header.tpl} (lines 147--148): the per-accelerator Config and State accessors.}]{../software/templates/versat_header.tpl}

Putting this together with the \texttt{SimpleExample} module of
Listing~\ref{lst:simple_example} (Section~\ref{sec:versat_spec}): its
generated \texttt{Config} struct has one member per instance (\texttt{a}
and \texttt{b}, each a \texttt{ConstConfig} exposing the \texttt{constant}
Config wire of Listing~\ref{lst:const_unit}), and its \texttt{State} struct
exposes \texttt{result.currentValue}, the State wire that \texttt{Reg}
(\texttt{hardware/src/units/Reg.v}) derives directly from its output
register:

\begin{lstlisting}[language=C,caption={Illustrative instantiation and run sequence for \texttt{SimpleExample}.},label={lst:inst_example}]
#include "versat_accel.h"

int main(){
  int base = ...; // Accelerator's memory-mapped base address, provided by
                   // the surrounding IOb-SoC integration (Section 2.4).
  versat_init(base);

  accelConfig->a.constant = 3;
  accelConfig->b.constant = 4;

  RunAccelerator(1); // Configuration reaches the datapath one run late
                      // (Section~\ref{sec:if_config_state}), so this first
                      // run still computes with the reset configuration.

  int stale = accelState->result.currentValue; // stale == 0

  RunAccelerator(1); // Second run: now computes with the values written above.

  int sum = accelState->result.currentValue; // sum == 7

  return 0;
}
\end{lstlisting}

The first run above reads back \texttt{0}, not the sum, because of the
shadow-register pipeline explained in Section~\ref{sec:if_config_state}:
software's write only reaches the datapath at the start of the
\emph{next} run. Section~\ref{sec:simple_example_software} walks through
this same accelerator's actual output, run by run.

\texttt{StartAccelerator}/\texttt{EndAccelerator} let software overlap the
CPU's own work with the accelerator's, as described in
Section~\ref{sec:if_config_state}. \texttt{RunAccelerator} above is simply
a convenience wrapper: \texttt{StartAccelerator} immediately followed by
\texttt{EndAccelerator}, for callers with nothing else to do while the
accelerator runs.

\subsubsection*{Config and State Addresses}

On the embedded target, where \texttt{iptr} is 4 bytes wide, \texttt{accelConfig} and \texttt{accelState} resolve to real, fixed byte addresses relative to \texttt{versat\_base}. Table~\ref{tab:simple_example_addrs} gives them for \texttt{SimpleExample}, read directly from the address decode logic Versat generated in \texttt{versat\_configurations.v} and \texttt{versat\_instance.v}. PC emulation does not use this address map at all: there, \texttt{accelConfig} and \texttt{accelState} point at plain host memory buffers instead, since there is no real memory-mapped interface to decode addresses against.

\begin{table}[H]
\centering
\begin{tabularx}{\textwidth}{|c|c|l|X|}
\hline
\rowcolor{iob-green}
\textbf{Address} & \textbf{Direction} & \textbf{Field} & \textbf{Notes} \\
\hline
\hline
0x00 & Read/Write & control register & write byte 0 to start a run; write byte 3, bit 31 for soft reset, bit 30 for signal loop; read bit 0 for \texttt{done} \\
\hline
0x04 & Write & \texttt{a.constant} & Config \\
\hline
0x04 & Read & \texttt{result.currentValue} & State; same address as \texttt{a.constant}, opposite direction \\
\hline
0x08 & Write & \texttt{b.constant} & Config \\
\hline
0x0C & Write & \texttt{result.disabled} & Config \\
\hline
0x10 & Write & \texttt{Delay0} & Config; internal, not exposed in \texttt{SimpleExampleConfig} \\
\hline
\end{tabularx}
\caption{Config and State byte addresses for \texttt{SimpleExample}, relative to \texttt{versat\_base}.}
\label{tab:simple_example_addrs}
\end{table}

Address \texttt{0x04} is read and written for two different purposes: a write there lands in the shadow copy of \texttt{a.constant}, subject to the one-run delay explained above, while a read there returns \texttt{result.currentValue} directly and combinationally, with no such delay on the read side; only when a Config write takes effect is delayed, never how promptly a State read reflects the accelerator's current output. Section~\ref{sec:if_mem_mapped} covers the separate, optional memory-mapped path into a unit's own native interface (reachable here at \texttt{versat\_base + 0x40}), which this example does not use.
